%gibt an: Papierformat, einseitiger Druck, Schriftgr��e
\documentclass[a4paper,oneside,titlepage,12pt]{article}
%-------------------------------------------------------------------
\usepackage[a4paper, top=2cm, footskip=0pt, headheight=0.8cm, headsep=0.6cm, lmargin=3cm, rmargin=2cm]{geometry}
\usepackage{graphicx}
\usepackage{helvet}
\usepackage{amsmath}
\usepackage{amsthm}
\usepackage{amssymb}
\usepackage{hyperref} 
\usepackage[right]{eurosym}
\usepackage[latin1]{inputenc}

%--------------------------------------------------------------------
\renewcommand{\baselinestretch}{1.2}
%damit unten Platz f�r die Seitennummer bleibt

\begin{document}
%--------------------------------------------------------------------
%Titelseite
\begin{titlepage}
	\includegraphics{grafiken/HTW-Logo.png}
	%includegraphics[width=.3\textwidth]{grafiken/HTW-Logo.png}
	\vspace*{3cm}
	\begin{center}
		\Huge{Beleg - Dokumentation\\} \vspace*{1cm}
		\Large{
			Modul: Software Engineering II\\
			Entwicklung eines Software-Teilssystems zur Verwaltung
			von Beleggruppendaten\\}
		Chrisitan Knothe\\Martin Tzschoppe\\Benjamin Reim\\Markus
		Noack\\Benjamin Herzog\\Christian Schwarz\\Felix Krautschuk\\ 
		\vspace*{1cm}
		
		\vspace*{2cm}
		\normalsize{
			Studiengang Informatik\\
		}
	\end{center}
	\vspace{2cm}
\begin{center}
\large{Sommersemester 2014}
\end{center}
	\vspace*{3cm}



\end{titlepage}
\setcounter{page}{2} \pagenumbering{arabic}

\thispagestyle{empty}\clearpage

%-----------------------------------------------------------------------------------------------------------------------------
%\rmfamily \pagestyle{fancy} \setcounter{secnumdepth}{4}
\newtheorem{satz}{Satz}
\newtheorem{lemma}[satz]{Lemma}
\newtheorem{folgerung}[satz]{Folgerung}
\theoremstyle{definition}
\newtheorem{definition}[satz]{Definition}
\numberwithin{equation}{section}
\renewcommand{\proofname}{Beweis}

\pagenumbering{roman}\setcounter{page}{3} \tableofcontents
\newcounter{roemisch} \setcounter{roemisch}{\value{page}}
\clearpage
\section{Verantwortlichkeiten}
Projektleiter: Christian Knothe\\
Analyse: Martin Tzschoppe\\
Entwurf: Benjamin Reim\\
Datenbank: Markus Noack\\
Implementierung: Benjamin Herzog\\
Tests: Christian Schwarz\\
Dokumentation/Protokollierung: Felix Krautschuk
\section{Gruppen-Besprechungen}
\subsection{1. Treffen}
Termin: 07.04.2014\\
Tagesordung:
\begin{itemize}
  \item Verantwortlichkeiten festlegen
  \item Wahl des zu bearbeitenden Projekts
  \item Einigung auf Kommunikationsverfahren und -infrastrukturen
  \item Diskussionen �ber das Vorgehensmodell (agile Softwareentwicklung,
  Wasserfallmodell,\ldots)\\
\end{itemize}
Resultat:
\begin{itemize}
  \item Verantwortlichkeiten: siehe Auflistung
  \item Thema: Verwaltung von Beleggruppendaten (Alternative 5)
  \item Vorgehensweise: Wasserfallmodell
\end{itemize}

\newpage
\subsection{2. Treffen}
Termin: 23.04.2014 in der Bibliothek B302a\\
Tagesordnung:
\begin{itemize}
  \item erste Ideen f�r die Analyse der Problemstellung zusammentragen
  \item Klarheit �ber Komplexit�t und Detailliertheit der Features des Systems
  \item Entwurfsm�glichkeiten diskutieren
  verschaffen\\
\end{itemize}
Probleme die diskutiert wurden, aber noch mit Frau Hauptmann zu besprechen
sind:\\Wie sollen die Studenten-Gruppen �berhaupt verwaltet werden:
\begin{itemize}
  \item 1. M�glichkeit: Dozent tr�gt die Mitglieder der Gruppe selbst in das
  System ein
  \begin{itemize}
  	\item Dozent hat einen guten �berblick �ber alle Gruppen und sieht sofort,
    welche Gruppen bereits vollst�ndig sind und wo noch Mitglieder fehlen und
    welche Studenten noch keiner Gruppe zugeordnet wurden
    \item aber: dadurch hat der Dozent am meisten Aufwand bei der Organisation
    der Gruppen
  \end{itemize}
  \item 2. M�glichkeit: Dozent erstellt f�r die Gruppen eine Art Grundger�st
  (Rahmen) und die Studenten tragen dann in diese zun�chst leeren Gruppen ein
  \begin{itemize}
    \item  Dozent bleibt beim Eintragen der Gruppen passiv, hat nur
    Kontrollfunktion
    \item m�sste das Eintragen dann online (in einer Webanwendung) geschehen??
    \item w�re schwieriger zu implementieren als 1. Variante
  \end{itemize}
  \item 3. M�glichkeit: potentielle Gruppenleiter melden sich bei Dozenten und
  werden eingetragen
  \begin{itemize}
  \item diese suchen sich ihre Mitglieder selbst
  \item es gibt somit keine leeren Gruppen
  \item Organisation der Gruppen geschieht ebenfalls durch Studenten selbst
  \end{itemize}
\end{itemize}
Filterm�glichkeiten??\\
wie komplex soll das System insgesamt werden?? Viele Features oder einfach
halten??\\
Was umfasst der Entwurf alles??\\
----------------------\\
Plan f�r n�chstes Treffen:\\
Termin: 30.04.2014\\
Benjamin Reim: Dialoge, eventuell erste Entw�rfe\\
Markus Noack: Archivierung, alles was zum Schluss aus der Datenbank in der
PDF-Datei stehen muss

\newpage
\subsection{3. Treffen}
Termin: 30.04.2014 13.30-15.00 in der Bibliothek B302a\\
Tagesordnung:
\begin{itemize}
  \item Diskussion bez�glich des Aufbaus und Einsatzes des Systems mit Frau
  Prof. Hauptmann
  \item bisherige Zwischenergebnisse der Analayse an neue Informationen anpassen
\end{itemize}
Resultate nach dem Gespr�ch mit Frau Prof. Hauptmann:\\
Das Ziel des Projektes ist eine kleine, schlanke, lokale, m�glichst
unter Windows laufende Anwendung zur Automatisierung der ANMELDUNG der
Beleggruppen OHNE Benotung. Diese Anwendung soll von Studenten von Hochschule
aus und NICHT VON ZU HAUSE aus bedient werden! Der Dozent soll ohne gro�e M�he
auch nach langer Zeit herausfinden k�nnen in welcher Gruppe und mit welcher
Verantwortlichkeit eine Person XY gearbeitet hat.\\
Der Dozent soll mithilfe von Filter-Funktionen Emails versenden k�nnen, also
z.B. an "alle Gruppenleiter", "alle Analytiker", "Gruppe case04 und Gruppe 07",
\ldots

Informationen, die im System sichtbar werden sollen:
\begin{itemize}
  \item insgesamt:
  \begin{itemize}
    \item Welche Beleg-Gruppen gibt es?
    \item Wer geh�rt zu welcher Gruppe?
  \end{itemize}
  \item f�r einzelne Gruppen im Detail
  \begin{itemize}
    \item Wer ist der Gruppenleiter?
    \item Welche Rollen haben die einzelnen Mitglieder?
  \end{itemize}
\end{itemize}

Wie verl�uft die Anmeldung einer Beleggruppe im System?
\begin{itemize}
  \item Dozent gibt in der Vorlesung bekannt dass die Anmeldung der Gruppen
  beginnen kann $\rightarrow$ f�r den Zeitraum der Anmeldung ist ein Start-Datum
  und ein End-Datum definiert
  \item ebenfalls in dieser Vorlesung wird der Erst-Login bekanntgegeben
  \item nachher (oder evtl bereits geschehen) finden sich die Gruppen zusammen
  \item Gruppe findet einen Gruppenleiter, welcher dann Ansprechpartner f�r
  Dozent darstellt
  \item Gruppe tr�gt sich in System ein und �ndert das Passwort (entweder kennen
  alle Mitglieder das neue Passwort oder nur der Gruppenleiter)
  \item es ist eine Mindest- und Maximalanzahl an Gruppenmitglieder definiert
  \item innerhalb des Anmelde-Zeitraumes kann Gruppe aus einem Pool ein zu
  bearbeitendes Thema ausw�hlen und dieses flexibel �ndern (mehrere Gruppen
  d�rfen das selbe Thema ausw�hlen)
  \item nach Beendigung der Belegarbeit (Ende des Semesters?) werden alle
  Beleg-Gruppendaten in einer PDF archiviert und die Datenbank bis auf die
  Themenauswahl ger�umt werden
\end{itemize}

Welche Informationen werden zur Anmeldung ben�tigt?
\begin{itemize}
  \item Gruppen-Identifikation (z.B. eine Gruppen-Nummer)
  \item Gruppenleiter
  \item ge�ndertes Passwort
  \item case-Nummer (Rahmenbedingungen f�r den Wertebereich!)
  \item pro Mitglied:
  \begin{itemize}
    \item Name
    \item Vorname
    \item s-Nummer
    \item Email-Adresse
    \item Veratnwortlichkeit (kann zum Zeitpunkt der Anmeldung noch
    freigelassen werden)
  \end{itemize}
\end{itemize}


\end{document}
